% !TeX root = main.tex

\subsection{SMT-LIB Standard for Unicode Theory of Strings}%

We begin with the input language of OSTRICH2\@.
It is based on SMT-LIB 2.6~\cite{SMTLIB} --- in particular, including support 
for the SMT-LIB Unicode theory and Linear Integer Arithmetic constraints
--- but additionally also supports some ``complex'' string functions including
transductions.
A formal description of the grammar is deferred to Appendix~\ref{app:grammar}, 
and here we illustrate the supported language features by example.

\subsubsection{SMT-LIB Constraints}

Below is a minimal SMT‑LIB script to show some core features.
The script uses the quantifier-free string theory with linear integer arithmetic (\verb+QF_SLIA+), contains three string variables \verb+x+, \verb+y+, and \verb+z+, and one integer variable \verb+l+.
It \verb+assert+s constraints on the variables that are explained below.
The final statements check the satisfiability of the constraints and produce a satisfying assignment.

\begin{lstlisting}[language=SMTLib,frame=single]
(set-logic QF_SLIA)

(declare-fun x () String)
(declare-fun y () String)
(declare-fun z () String)
(declare-fun k () Int)

(assert (= (str.len x) (+ k 1)))
(assert (= x (str.++ y z)))
(assert (str.in_re x (re.+ (re.union (str.to_re "a") (str.to_re "b")))))

(check-sat)
(get-model)
\end{lstlisting}

Constraints are written in a Lisp-style prefix notation.
The first \verb+assert+ requires that the length of \verb+x+ is equal to $\mathtt{k} + 1$.
The second requires that \verb+x+ is equal to the concatenation of \verb+y+ and \verb+z+.
The third requires that \verb+x+ belongs to the regular language
$(\{\mathtt{a}\} \cup \{\mathtt{b}\})^+$,
represents a non-empty sequence of $\mathtt{a}$ and $\mathtt{b}$ characters.

Standard Boolean connectives are supported.
The example below requires \verb+x+ to either be the concatenation of \verb+y+ and \verb+z+ or the concatenation of \verb+z+ and \verb+y+.
\begin{lstlisting}[language=SMTLib,frame=single]
(assert (or (= x (str.++ y z))) (= x (str.++ z y))))
\end{lstlisting}

\subsection{Extensions Beyond the SMT-LIB Standard}

OSTRICH2 provides additional features to model string constraints in practical applications.
A simple example is the \verb+str.reverse+ function.
Below, we introduce the transducer, regular expression, and automata extensions.

\subsubsection{Transducer-based Operations}

OSTRICH2 supports transducers and prioritised finite-state transducers~\cite{BDM14,BM17,popl22}.

A transducer (a.k.a. a \emph{rational transducer}) takes a single string as 
input and produces a single output string.
They enable operations such as HTML encoding (e.g.\ replacing \verb+&+ with \verb+&amp;+) to be defined.
Transducers are written as recursive functions with two arguments: the input and the output.
The example below shows a \verb+toUpper+ transduction.
The base case (line~4) applies when both strings are empty.
The recursive case (lines~5--11) asserts three conditions.
First, both arguments are non-empty.
Second, the head of the output (\verb+y+) is equal to the upper case version of the head of the input (\verb+x+).
Third, the tails of the two arguments recursively satisfy \verb+toUpper+.
To obtain the upper case version of a character, arithmetic is
performed on the character codes. In general, a transducer definition
can
contain multiple mutually recursive functions. The option
\verb!:parse-transducers! instructs OSTRICH2 to translate the recursive
function definition to an internal transducer representation and apply
an automata-based decision procedure to solve the constraints in lines~16--17.

\begin{lstlisting}[language=SMTLib,frame=single]
(set-option :parse-transducers true)

(define-fun-rec toUpper ((x String) (y String)) Bool
(or (and (= x "") (= y ""))
  (and (not (= x "")) (not (= y ""))
       (= (char.code (str.head y))
          (ite (and (<= 97 (char.code (str.head x)))
                    (<= (char.code (str.head x)) 122))
               (- (char.code (str.head x)) 32)
               (char.code (str.head x))))
       (toUpper (str.tail x) (str.tail y)))))

(declare-fun x () String)
(declare-fun y () String)

(assert (= x "Hello World"))
(assert (toUpper x y))
\end{lstlisting}

Prioritised transducers additionally allow certain transitions to take
precedence over other transitions.
That is, a non-deterministic transduction may only succeed on one path if another cannot succeed.
Such transducers are used internally in regular expression handling,
described below, but can also be formulated using the
\verb!define-fun-rec! notation (not shown here). We discuss these below after we
cover extended regular expressions.

\subsubsection{Extended Regular Expression Support}

The SMT-LIB standard defines a minimal set of regular expression
constructs,
all of which are supported by OSTRICH2.
In addition, OSTRICH2 can also parse ECMAScript regular expressions,
following the 2020 version~\cite{ECMAScript11}, and supports regular
expression features such as look-arounds that cannot directly
expressed in SMT-LIB.
A simple usage of ECMAScript regular expressions is shown below, using the \verb+re.from_ecma2020+ function.
The expression matches either sequences of characters or sequences of digits.
\begin{lstlisting}[language=SMTLib,frame=single]
(assert (str.in_re w (re.from_ecma2020 "[a-z]*|[0-9]*")))
\end{lstlisting}
To enable full support of ECMAScript character escaping, OSTRICH2 also supports
single-quoted strings, in which the normal SMT-LIB character escaping
is disabled and strings are passed unmodified to the
regular expression parser:
\begin{lstlisting}[language=SMTLib,frame=single]
(assert (str.in_re w (re.from_ecma2020 '\x24[0-9]+')))
\end{lstlisting}

One notable extension is the use of capture groups and references in replace and extract functions.
The following example shows the replacement of any substring matching \verb+src='.*'+ with the text appearing between the single quotation marks.
For example
%
\verb+<img src='image.png'>+ would become
%
\verb+<img image.png>+.
%
We use the \verb+re.from_ecma2020+ function for convenience;
direct functions for regular expression operators---such as \verb+re.capture+---are also available.
\begin{lstlisting}[language=SMTLib,frame=single]
(assert (= x
           (str.replace_cg_all
              y
              (re.from_ecma2020 "src='(.*)'")
              (_ re.reference 1))))
\end{lstlisting}
%\todo[inline]{MH: do we need (\_ re.reference 1) or is there a
%  fromECMA alternative? PR: not at the moment, but would be easy to add}

There are several features to note about this example.
The first argument of \verb+str.replace_cg_all+ is the string to be searched.
The second is the regular expression that captures part of the string using the parenthesis notation \verb+(+\ldots\verb+)+.
The final argument is the \emph{replacement pattern} that can be a concatenation of string literals and references to captured text.
In this case, only the contents of the first (and only) capture group are used.

Importantly, when matching \verb+src='(.*)'+, precedence rules must be respected.
The \verb+*+ operator is \emph{greedy} in ECMAScript (and many other regular expression languages), which means that it should match as many characters as possible before the remainder of the match begins.
This can be a cause of errors when developers write regular expressions and a challenge for symbolic execution tools using a string constraint solver~\cite{LMK19}.
For example, if \verb+y+ contained the value \verb+src='a' src='b'+, then there should only be one match of \verb+src='(.*)'+ rather than two.
In the match, the captured value is \verb+a' src='b+,
that is, the longest string surrounded by \verb+'+ symbols.
This semantics is respected by OSTRICH2 through the use of \emph{prioritised} transducers~\cite{BDM14,BM17,popl22}.
In a prioritised transducer, steps of the transduction can be given priority over the alternatives.
When matching against \verb+src='a' src='b'+ there is a choice whether the match of \verb+(.*)+ should stop at the second \verb+'+ or continue.
Because \verb+*+ is greedy, the match of \verb+(.*)+ has priority and
matching can only stop if continuing would fail. When replacing the
pattern
%
\verb+src='(.*)'+ with \verb+src='(.*?)'+,
%
using a \emph{lazy} quantifier~\verb+*?+, instead the
shortest string will be matched.

\subsubsection{Automata Representations}

OSTRICH2 also directly supports finite-state automata, which are often
more convenient than regular expressions when integrating an SMT
solver in applications.
The function \verb|re.from_automaton| parses a finite‑state automaton given as a string.
In the example below the automaton has
\begin{itemize}
    \item an initial state \texttt{s0},
    \item an accepting state \texttt{s1},
    \item a transition $\texttt{s0 -> s1 [0,100];}$ (accepting character codes 0–100),
    \item a loop $\texttt{s1 -> s1 [0,65535];}$ (accepting any Unicode character).
\end{itemize}
Then \verb|str.in_re| can be used to test whether the value of an expression belongs to the language of that automaton.

\begin{lstlisting}[language=SMTLib,frame=single]
(assert (str.in_re x
       (re.from_automaton "automaton value_0 {init s0; s0 -> s1 [0, 100]; s1 ->s1[0,65535]; accepting s1;};")))
\end{lstlisting}






